\paragraph{Theorem T5: Finite-Horizon Certificate Validity.}
Let \(X_t\) be the observation-consistent state set at issue time \(t\), let
\(\tilde a_{t:t+h}\) be the certified action or fallback sequence, and let
\(S\) be the domain safe-state set. Assume the forward reachable tube satisfies
\[
  \mathcal R_{t+j}(X_t,\tilde a_{t:t+j}) \subseteq S
  \qquad j=0,\ldots,h .
\]
Then every release governed by the certificate is safe for \(h\) steps unless a
certificate-invalidating event occurs.

\paragraph{Proof.}
At \(j=0\), all states consistent with the issued certificate lie in
\(\mathcal R_t(X_t)\subseteq S\). Suppose the claim holds through step \(j<h\).
By definition of the reachable tube, every state reachable after applying the
certified prefix \(\tilde a_{t:t+j+1}\) lies in
\(\mathcal R_{t+j+1}(X_t,\tilde a_{t:t+j+1})\). The containment assumption places
that set inside \(S\). Induction proves safety for all \(j=0,\ldots,h\). If a
fresh observation contradicts \(X_t\), the model, or the action sequence, the
certificate is outside the theorem hypothesis and must be renewed, repaired, or
expired by the runtime.
